
\subsection*{2.3}
\phantomsection
\label{sga3:VIIB:2.3}
\addcontentsline{toc}{subsection}{2.3. The covariant algebra}

Let \(G\) now be an arbitrary formal \(k\)-group,\footnote{
\label{sga3:VIIB:nde:85}That is, \(G\) is not necessarily
topologically flat over \(k\).} with affine algebra \(A\). We
continue to write
\[
  \mathbf H(G)=\mathbf V_k^f(A)
\]
for the \(\mathbf O_k\)-module dual to \(A\), and denote by
\[
  \phi_G:
  \mathbf H(G)\otimes_k\mathbf H(G)
  \longrightarrow
  \mathbf H(G\times G)
\]
the functorial homomorphism induced, for every object \(C\) of
\(\operatorname{Alf}/k\), by the natural map of \SGARefLink{sga3:VIIB:0.3.6}{0.3.6}:
\[
\begin{aligned}
  &(A\widehat\otimes_kC)^\dagger
  \otimes_C
  (A\widehat\otimes_kC)^\dagger
  \\
  &\hspace{3em}\longrightarrow
  \left(
    (A\widehat\otimes_kC)
    \widehat\otimes_C
    (A\widehat\otimes_kC)
  \right)^\dagger.
\end{aligned}
\]

If
\[
  m:G\times G\longrightarrow G
\]
is the multiplication of \(G\), the composite
\[
  \mathbf H(G)\otimes_k\mathbf H(G)
  \xrightarrow{\ \phi_G\ }
  \mathbf H(G\times G)
  \xrightarrow{\ \mathbf H(m)\ }
  \mathbf H(G)
\]
endows \(\mathbf H(G)\) with an algebra structure over
\(\mathbf O_k\). For every
\(C\in\operatorname{Alf}/k\), the unit element of
\[
  \mathbf H(G)(C)
  =
  (A\widehat\otimes_kC)^\dagger
\]
is the augmentation of \(A\widehat\otimes_kC\); cf.\ \SGARefLink{sga3:VIIA:2.1}{2.1}.
\footnote{\label{sga3:VIIB:nde:86}The editors changed the order of
the sentences that follow.}

If \(G\) is not topologically flat over \(k\), the morphism
\(\phi_G\) need not be an isomorphism. Consequently, the morphism
\[
  \delta_G:
  \mathbf H(G)\longrightarrow\mathbf H(G\times G)
\]
induced by the diagonal \(x\mapsto(x,x)\) of \(G\) need not factor
through
\[
  \mathbf H(G)\otimes_k\mathbf H(G).
\]
Thus \(\mathbf H(G)\) need not be a bialgebra over
\(\mathbf O_k\). For this reason, in the general case we simply call
\(\mathbf H(G)\) the \emph{covariant algebra} of the formal group
\(G\).

When \(G\) is topologically flat over \(k\), of course, \(\phi_G\) is
an isomorphism, and one recovers the \(\mathbf O_k\)-bialgebra
structure on \(\mathbf H(G)\) defined in \SGARefLink{sga3:VIIB:2.2.1}{2.2.1}.
